% The short CV's look: the site's (DESIGN.md). A dark sea at night and four lights on it.
% lualatex. The three faces are the site's: Bricolage Grotesque, Newsreader, JetBrains Mono;
% cv/build.sh copies them from the flake's fonts into \cvfonts. \cvtitle comes from vars.tex.
\usepackage[a4paper,margin=15mm,top=12mm,bottom=20mm,footskip=11mm]{geometry}
\usepackage{fontspec}
\setmainfont{Newsreader.ttf}[Path=\cvfonts/, RawFeature={axis={wght=400,opsz=14}},
  BoldFont=Newsreader.ttf, BoldFeatures={RawFeature={axis={wght=650,opsz=14}}},
  ItalicFont=Newsreader-Italic.ttf, ItalicFeatures={RawFeature={axis={wght=400,opsz=14}}}]
\newfontfamily\display{Bricolage.ttf}[Path=\cvfonts/, RawFeature={axis={wght=800,opsz=96,wdth=100}}]
\newfontfamily\displaysemi{Bricolage.ttf}[Path=\cvfonts/, RawFeature={axis={wght=600,opsz=24,wdth=100}}]
\newfontfamily\mono{JetBrainsMono-Regular.ttf}[Path=\cvfonts/, Scale=0.82]
\setmonofont{JetBrainsMono-Regular.ttf}[Path=\cvfonts/, Scale=0.82]
\usepackage{xcolor}
% The ground and the four lights, the same tokens as site/static/style.css.
\definecolor{night}{HTML}{05090E}
\definecolor{rule}{HTML}{1C2B38}
\definecolor{ink}{HTML}{E7F2EF}
\definecolor{muted}{HTML}{8FA3AD}
\definecolor{water}{HTML}{36B7FF}
\definecolor{ideas}{HTML}{3DFF96}
\definecolor{ouroboric}{HTML}{FF8A1F}
\definecolor{closed}{HTML}{FF4D62}
\pagecolor{night}\color{ink}
\usepackage{graphicx}
\usepackage{multicol}
\usepackage{luacolor}
\usepackage{lua-ul}
\usepackage{enumitem}
\setlist[itemize]{nosep, leftmargin=1.2em}
\usepackage{hyperref}
\hypersetup{hidelinks, pdftitle={\cvtitle}, pdfauthor={Matheus Hoffmann}}
\pagestyle{empty}
\setlength{\parindent}{0pt}\setlength{\parskip}{4pt}
\linespread{1.12}\raggedright
\setlength{\columnsep}{9mm}
% A link is the words, with a line under them: a light is for meaning, not for "press here".
\newcommand{\cvlink}[1]{{\underLine[color=muted,top=-1.6pt,bottom=-2pt]{#1}}}
\newcommand{\cvdata}{\mono\color{muted}\small}
\newcommand{\cvdoi}[1]{{\mono\footnotesize\color{muted}\underLine[color=rule,top=-1.6pt,bottom=-2pt]{#1}}}
% The name is green and the epithet is blue (DESIGN.md, rule 3).
\newcommand{\cvname}[1]{{\display\fontsize{40}{40}\selectfont\color{ideas}#1\par}\vspace{-3pt}}
\newenvironment{cvepithet}{\displaysemi\fontsize{15}{17}\selectfont\color{water}}{\par\vspace{5pt}}
\newenvironment{cvlede}{\fontsize{12}{17.5}\selectfont}{\par\vspace{4pt}}
\newcommand{\cvpill}[1]{\mbox{\setlength{\fboxsep}{0pt}\cvpillbox{\mono\footnotesize\color{ink}#1}}}
% a pill: a round outline in the colour of a rule
\usepackage{tikz}
\newcommand{\cvpillbox}[1]{\tikz[baseline=(p.base)]\node[draw=muted, line width=0.5pt, rounded corners=7pt,
  inner xsep=7pt, inner ysep=4pt, text=ink] (p) {#1};}
\newenvironment{cvcontact}{\setlength{\parskip}{0pt}\spaceskip=5pt}{\par\vspace{2pt}}
\newcommand{\cvsection}[1]{\vspace{9pt}{\displaysemi\fontsize{15}{17}\selectfont\color{white}#1\par}%
  \vspace{-5pt}{\color{rule}\rule{\linewidth}{0.6pt}}\par\vspace{0pt}}
\newlength{\cvwhen}\setlength{\cvwhen}{25mm}
\newcommand{\cvrow}[2]{\makebox[\cvwhen][l]{\cvdata #1}\parbox[t]{\dimexpr\linewidth-\cvwhen}{\raggedright #2}\par\vspace{3pt}}
% A work: a light on its left edge, not a card (DESIGN.md, shapes).
\newsavebox{\cvworkbox}
\newenvironment{cvwork}[1]{\def\cvlight{#1}\begin{lrbox}{\cvworkbox}\begin{minipage}[t]{\dimexpr\linewidth-10pt}%
  \setlength{\parskip}{2pt}\raggedright}{\end{minipage}\end{lrbox}%
  \par\noindent{\color{\cvlight}\vrule width 2pt}\hspace{8pt}\usebox{\cvworkbox}\par\vspace{7pt}}
\newenvironment{cvworks}{\begin{multicols}{2}}{\end{multicols}\vspace{-8pt}}
% Two short sections side by side.
\newenvironment{cvpair}{\vspace{2pt}\begin{multicols}{2}\raggedcolumns}{\end{multicols}}
% The tip is the foot of the page: it is set in a box and the page's foot shows the box.
\newsavebox{\cvtipbox}
\makeatletter
\newenvironment{cvtip}{\begin{lrbox}{\cvtipbox}\begin{minipage}[b]{\linewidth}\setlength{\parskip}{0pt}%
  {\color{rule}\rule{\linewidth}{0.6pt}}\par\vspace{2pt}\footnotesize\color{muted}\raggedright}%
  {\end{minipage}\end{lrbox}\global\setbox\cvtipbox\copy\cvtipbox\gdef\@oddfoot{\usebox{\cvtipbox}}}
\makeatother
